\section{Kiến trúc tổng thể}

Kiến trúc tổng thể được thể hiện trong Hình~\ref{fig:platform-architecture}.
Người dùng làm việc qua Web Dashboard; mọi yêu cầu đi qua reverse proxy Nginx
tới Cloud API, WebSocket hoặc AI Assistant. Cloud Server là một ứng dụng FastAPI
duy nhất phục vụ đồng thời REST \texttt{/api/v1}, WebSocket \texttt{/api/v1/ws},
MCP server \texttt{/mcp} và bộ tiêu thụ MQTT (embedded worker). Trạng thái bền
vững nằm trong PostgreSQL, telemetry ngắn hạn nằm trong Redis, còn kênh Cloud--Edge
là MQTT broker. Model được lưu trên Hugging Face Hub; image runtime nằm trên
GitHub Container Registry (GHCR).

\begin{figure}[H]
\centering
\resizebox{\textwidth}{!}{%
\begin{tikzpicture}[font=\small,
  svc/.style={draw, rounded corners, align=center, minimum height=1cm, text width=2.8cm, fill=blue!5},
  store/.style={draw, align=center, minimum height=0.9cm, text width=2.4cm, fill=orange!8},
  ext/.style={draw, dashed, rounded corners, align=center, minimum height=0.9cm, text width=2.6cm, fill=gray!6},
  edge/.style={draw, rounded corners, align=center, minimum height=1.6cm, text width=3.6cm, fill=purple!5},
  lbl/.style={font=\scriptsize, fill=white, inner sep=1pt}]
  % Hàng người dùng và Cloud
  \node[svc, fill=green!6] (user) at (0,1.6) {Người vận hành\\(trình duyệt)};
  \node[svc] (web) at (0,-0.6) {Web Dashboard\\React + Nginx};
  \node[svc, text width=4.3cm, minimum height=2.9cm] (cloud) at (5.2,-0.6)
      {\textbf{Cloud Server} (FastAPI)\\REST /api/v1\\WebSocket /api/v1/ws\\MCP /mcp\\Embedded MQTT worker};
  \node[svc] (assistant) at (5.2,2.4) {AI Assistant\\LangGraph + LLM};
  % Kho dữ liệu
  \node[store] (pg) at (10.0,0.8) {PostgreSQL 16};
  \node[store] (redis) at (10.0,-0.6) {Redis 7};
  \node[store] (mqtt) at (10.0,-2.0) {MQTT Broker\\Mosquitto 2};
  % Dịch vụ ngoài
  \node[ext, text width=3.0cm] (ghcr) at (13.9,1.0) {GHCR (image runtime)};
  \node[ext, text width=3.0cm] (hf) at (13.9,-0.8) {Hugging Face Hub (kho model)};
  % Edge
  \node[edge] (edge1) at (6.0,-5.4) {\textbf{Edge Node 1}\\Edge Agent\\AI Runtime (container)\\Camera};
  \node[edge] (edge2) at (11.0,-5.4) {\textbf{Edge Node N}\\Edge Agent\\AI Runtime (container)\\Camera};

  \draw[flowarrow] (user) -- (web);
  \draw[flowarrow] (web) -- node[lbl, below]{HTTPS/WS} (cloud);
  \draw[flowarrow] (user.east) -- node[lbl, above]{SSE} (assistant.west);
  \draw[flowarrow] (assistant) -- node[lbl, right]{MCP} (cloud);
  \draw[flowarrow] (cloud.east |- pg) -- (pg);
  \draw[flowarrow] (cloud.east |- redis) -- (redis);
  \draw[flowarrow, <->] (cloud.east |- mqtt) -- (mqtt);
  \draw[flowarrow, <->] (mqtt.south) -- ++(0,-1.0) -| node[lbl, pos=0.25, above]{MQTT} (edge1.north);
  \draw[flowarrow, <->] (mqtt.south) -- ++(0,-1.0) -| (edge2.north);
  \draw[flowarrow, dashed] (hf.south) -- (13.9,-3.9) -- node[lbl, above]{tải model} (12.55,-3.9) -- (12.55,-4.6);
  \draw[flowarrow, dashed] (ghcr.east) -- (16.0,1.0) -- (16.0,-6.9) -- node[lbl, below, pos=0.5]{pull image runtime} (6.0,-6.9) -- (edge1.south);
\end{tikzpicture}}
\caption{Kiến trúc tổng thể nền tảng Edge AI}
\label{fig:platform-architecture}
\end{figure}

\subsection{Control Plane và Data Plane}

\textbf{Control Plane} (Cloud) chịu trách nhiệm quản lý tài nguyên, lưu desired
state của từng runtime dưới dạng Deployment, phát lệnh, tổng hợp trạng thái và
cung cấp dữ liệu giám sát. \textbf{Data Plane} (Edge) chịu trách nhiệm xử lý dữ
liệu tại chỗ: thu nhận khung hình, suy luận, quản lý vòng đời container và thu
thập telemetry. Hai mặt phẳng chỉ trao đổi qua các thông điệp đã định nghĩa trong
gói hợp đồng dùng chung \texttt{src/shared/contracts}.

Luồng điều khiển đi theo chiều \textbf{Client $\rightarrow$ Cloud $\rightarrow$
Edge Agent $\rightarrow$ AI Runtime}; luồng telemetry đi theo chiều ngược lại
\textbf{AI Runtime $\rightarrow$ Edge Agent $\rightarrow$ Cloud}. AI Runtime không
kết nối trực tiếp tới broker mà gửi chỉ số qua HTTP tới API cục bộ của Agent; nhờ
đó chỉ Agent giữ credential MQTT và Cloud chỉ cần tin cậy một thực thể trên mỗi node.

\subsection{Giao thức MQTT và cấu trúc topic}

MQTT \cite{mqtt5} được chọn làm kênh Cloud--Edge vì nhẹ, hỗ trợ kết nối từ sau
NAT (Edge chủ động kết nối ra ngoài), có cơ chế retained message và phù hợp với
đường truyền không ổn định. Topic được tổ chức theo phân cấp tài nguyên như
Bảng~\ref{tab:mqtt-topics}.

\begin{table}[H]
\centering
\caption{Cấu trúc topic MQTT giữa Cloud và Edge}
\label{tab:mqtt-topics}
\small
\begin{tabularx}{\textwidth}{L{6.6cm}L{2.5cm}Y}
\toprule
\textbf{Topic} & \textbf{Chiều} & \textbf{Nội dung}\\
\midrule
\path{edge/{cluster}/{node}/command/{type}} & Cloud$\rightarrow$Edge & Lệnh \texttt{deploy}, \texttt{update}, \texttt{start}, \texttt{stop}, \texttt{restart}, \texttt{remove}\\
\path{edge/{cluster}/{node}/state} & Edge$\rightarrow$Cloud & NodeState (retained), đồng thời là heartbeat\\
\path{edge/{cluster}/{node}/telemetry/{kind}} & Edge$\rightarrow$Cloud & \texttt{system}, \texttt{gpu}, \texttt{runtime}\\
\path{edge/{cluster}/{node}/event/{kind}} & Edge$\rightarrow$Cloud & \texttt{runtime}, \texttt{error}, \texttt{deployment}\\
\bottomrule
\end{tabularx}
\end{table}

Cloud đăng ký các wildcard \path{edge/+/+/state/#}, \path{edge/+/+/telemetry/+}
và \path{edge/+/+/event/+}; mỗi Agent chỉ đăng ký \path{edge/{cluster}/{node}/command/+}
của chính nó. Lớp \texttt{EdgeTopics} vừa sinh vừa phân tích topic, được cả Cloud
và Edge dùng chung.

\subsection{Tổ chức mã nguồn theo lớp}

Mã nguồn backend là một monorepo Python 3.12, tách phần ứng dụng khởi chạy
(\texttt{apps/}) khỏi phần nghiệp vụ (\texttt{src/}). Trong mỗi bounded context
(\texttt{cloud}, \texttt{edge}, \texttt{assistant}), mã được chia thành bốn lớp:

\begin{itemize}
  \item \textbf{domain}: entity, value object và interface repository, không phụ
        thuộc framework;
  \item \textbf{application}: service thực hiện use case, handler tiêu thụ thông
        điệp và \emph{port} (interface tới hạ tầng như container engine, camera,
        detector);
  \item \textbf{infrastructure}: \emph{adapter} hiện thực port và repository
        (PostgreSQL, Redis, Docker CLI, OpenCV, Ultralytics, MQTT);
  \item \textbf{container}: lớp lắp ráp phụ thuộc, chọn adapter theo cấu hình
        (ví dụ \texttt{storage\_backend=memory} để chạy không cần database).
\end{itemize}

Cách tổ chức port--adapter cho phép thay \texttt{DockerRuntimeDriver} bằng
\texttt{FakeRuntimeDriver}, hay \texttt{OpenCvFrameSource} bằng
\texttt{SyntheticFrameSource} khi kiểm thử, mà không sửa service nghiệp vụ.
Ba gói nội bộ dùng lại giữa các dịch vụ: \texttt{hopquangdo-messaging} (trừu
tượng PubSub trên MQTT), \texttt{hopquangdo-http} (xác thực API key, health
router, xử lý lỗi) và \texttt{hopquangdo-observability} (cấu hình log).
